% Generated draft: requires mathematical and source review.
\documentclass[11pt]{article}
\usepackage[margin=25mm]{geometry}
\usepackage{fontspec}
\setmainfont{Latin Modern Roman}
\usepackage{amsmath,amssymb,mathtools,unicode-math}
\setmathfont{Latin Modern Math}
\usepackage{enumitem,needspace,xurl,hyperref}
\setlength{\parindent}{0pt}
\setlength{\parskip}{5pt}
\setlength{\emergencystretch}{3em}
\newcommand{\N}{\mathbb{N}}
\newcommand{\Z}{\mathbb{Z}}
\newcommand{\Q}{\mathbb{Q}}
\newcommand{\R}{\mathbb{R}}
\newcommand{\C}{\mathbb{C}}
\newcommand{\F}{\mathbb{F}}
\newcommand{\A}{\mathbb{A}}
\newcommand{\PP}{\mathbb{P}}
\newcommand{\E}{\mathbb{E}}
\newcommand{\eps}{\varepsilon}
\newcommand{\dd}{\,\mathrm{d}}
\begin{document}
\section*{UnsolvedMath ID: 30006422}
\textbf{Source problem code:} catalogue entry
\subsection{Definitions and mathematical statement}
A complex \(C^*\)-algebra \(A\) is called a \(\mathrm{SAW}^*\)-algebra if, whenever \(a,b\in A\) are positive and \(ab=0\), there is a positive contraction \(e\in A\) such that \(ea=a\) and \(eb=0\). All tensor products below are minimal \(C^*\)-tensor products, denoted by \(\otimes_{\min}\), as explicitly stipulated in the supplied original passage. Isomorphisms mean \(*\)-isomorphisms. Two \(C^*\)-algebras \(A\) and \(B\) are stably isomorphic if
\[
A\otimes_{\min}\mathcal K(\ell^2(\N))\cong B\otimes_{\min}\mathcal K(\ell^2(\N)),
\]
where \(\mathcal K(\ell^2(\N))\) is the algebra of compact operators on \(\ell^2(\N)\).

For the finite, nonempty-product reading of the supplied question, let \(m,n\) be positive integers and let \(A_i\), \(0\leq i<m\), and \(B_j\), \(0\leq j<n\), be \(\mathrm{SAW}^*\)-algebras satisfying
\[
\bigotimes_{i=0}^{m-1}{}_{\min} A_i
\cong
\bigotimes_{j=0}^{n-1}{}_{\min} B_j.
\]
The question is whether this implies \(m=n\), and, if so, whether there exists a permutation \(\pi\) of \(\{0,\ldots,m-1\}\) such that
\[
A_i\otimes_{\min}\mathcal K(\ell^2(\N))
\cong
B_{\pi(i)}\otimes_{\min}\mathcal K(\ell^2(\N))
\qquad(0\leq i<m).
\]
The finite, positive-integer convention for \(m,n\) makes this reading explicit; the source does not specify their range or conventions for empty products. More substantially, the supplied formulation imposes no nonzero, infinite-dimensional, or other nontriviality hypothesis on the factors. Such restrictions must not be silently added, and the intended hypotheses require review before this is regarded as a precise uniqueness conjecture.

The record also asks the following concrete question. Let \(H\) be a separable, infinite-dimensional complex Hilbert space, let \(\mathcal B(H)\) and \(\mathcal K(H)\) denote its bounded and compact operator algebras, and define the Calkin algebra by
\[
\mathcal Q(H)=\mathcal B(H)/\mathcal K(H).
\]
Does there exist an isomorphism
\[
\mathcal Q(H)\otimes_{\min}\mathcal Q(H)
\cong
\mathcal Q(H)\otimes_{\min}\mathcal Q(H)\otimes_{\min}\mathcal Q(H)?
\]
The supplied passage attributes the motivating question \(\mathcal Q(H)\cong\mathcal Q(H)\otimes_{\min}\mathcal Q(H)\) to Simon Wasserman. The supplied dated assessment attributes to Kania a result excluding decomposition of the Calkin algebra into two infinite-dimensional \(C^*\)-algebra factors and describes the remaining factorization and uniqueness questions as open. These are supplied literature claims, not independently verified conclusions of this draft. References in the record to norm ultrapowers and tracial ultraproducts do not include enough data to formulate additional questions.
\subsection{Short English statement}
The record asks whether a finite tensor product of \(\mathrm{SAW}^*\)-algebras determines its number of factors and their stable isomorphism classes, up to permutation. It also asks whether the minimal tensor square and tensor cube of the Calkin algebra are isomorphic. The general question needs review of missing restrictions on the factors and index ranges. The supplied status is partially solved; its dated literature assessment has not been independently verified.
\subsection{Sources}
\begin{itemize}
\item Supplied local UnsolvedMath registry record, permanent ID \(30006422\), source problem code \texttt{catalogue entry}, entitled ``Tensor-Factorization Uniqueness for \(\mathrm{SAW}^*\)-Algebras,'' in the collection ``Oberwolfach Reports Open Problems.'' The record gives proposed year 2025 and no named proposer.
\item Supplied report citation: \emph{\(C^*\)-Algebras (2026)}, Oberwolfach Reports, pp. 1887--1948, DOI \texttt{10.4171/owr/2025/35}: \url{https://doi.org/10.4171/owr/2025/35}. This is also the supplied literature source labelled ``Oberwolfach Report 35.'' The citation's parenthetical year 2026 and the DOI's year component 2025 are retained without reconciliation.
\item Supplied literature source labelled ``Kania tensor decomposition theorem'': \url{https://doi.org/10.1093/qmath/hav027}. The supplied assessment attributes to Kania the exclusion of a decomposition of the Calkin algebra into two infinite-dimensional \(C^*\)-algebras and describes this as settling Wassermann's original decomposition question. The original passage spells the motivating questioner's name ``Simon Wasserman''; both supplied spellings are recorded without independent identification.
\item Literature assessment checked at \texttt{2026-08-22}: the supplied curation labels the problem ``partially solved'' and describes further ultrapower factorization and uniqueness questions as open. This is supplied dated curation, not independently verified. The background expressly calls for independent verification of time-sensitive status and cited claims.
\item The supplied registry labels the statement \texttt{corrected\_verified} and says that the official report was checked and its notation repaired. This is a supplied verification claim; no source visit or independent mathematical verification is asserted here.
\end{itemize}
\end{document}
